\section{Related Work}
\label{sec:related}

\paragraph{Structured pruning.}
Network pruning reduces model size by removing parameters or entire structural units. \emph{Unstructured pruning}~\citep{han2015learning} zeroes individual weights but yields sparse tensors that require specialized hardware. \emph{Structured pruning} removes entire filters~\citep{li2017pruning}, channels~\citep{liu2017learning}, or attention heads~\citep{chen2021chasing,yu2022width}, producing dense sub-networks that run efficiently on standard accelerators. Our work falls into the structured category, pruning entire neurons within MoE expert feed-forward layers.

\paragraph{Sparsity scheduling.}
\citet{zhu2017prune} introduced a polynomial schedule for gradually increasing sparsity during training, showing that the \emph{trajectory} of pruning matters as much as the final sparsity level. Their cubic schedule $s(t) = s_f(1-(1-t)^3)$ starts slowly and accelerates, giving the network time to identify which weights are important before committing to aggressive pruning. While originally applied to unstructured magnitude pruning, we adapt this schedule to \emph{structured} pruning of elastic neuron widths, using mask logits rather than weight magnitudes as the pruning criterion.

\paragraph{Lagrangian and regularization-based pruning.}
MorphNet~\citep{gordon2018morphnet} applies group-lasso regularization guided by a FLOPs objective. Budget-aware regularization~\citep{lemaire2019structured} introduces Lagrangian multipliers that adaptively control resource usage. $L_0$ regularization~\citep{louizos2018learning} learns gates using the hard-concrete distribution. Variational dropout~\citep{molchanov2017variational} prunes via learned dropout rates. More recently, MaskLLM~\citep{fang2024maskllm} proposed learnable semi-structured sparsity patterns for large language models using Gumbel-softmax sampling over mask distributions, demonstrating that differentiable mask learning can scale to billion-parameter models. These methods share a common assumption: that the optimization signal from the resource penalty is strong enough to guide per-parameter pruning decisions. We show this assumption breaks down at per-neuron granularity in MoE architectures, motivating our schedule-based alternative.

\paragraph{Mixture-of-Experts.}
The MoE paradigm~\citep{shazeer2017outrageously} routes tokens to specialized expert sub-networks, enabling model capacity to grow without proportional FLOPs increase. Switch Transformers~\citep{fedus2022switch} simplified routing to top-1 selection, while GShard~\citep{lepikhin2021gshard} scaled MoE to production settings. Vision MoE~\citep{riquelme2021scaling} brought sparse experts to image recognition. Recent work has explored expert collapse~\citep{chi2022representation} and expert pruning~\citep{kim2023mixture}, removing entire experts rather than neurons within experts. \citet{lu2024expertslimming} showed that not all experts contribute equally and proposed task-aware expert skipping for MoE LLMs. MegaBlocks~\citep{gale2023megablocks} introduced efficient GPU kernels for sparse MoE training, addressing the routing overhead that limits wall-clock gains at small scale. KaleidoNet takes the complementary approach of pruning \emph{within} each expert at the neuron level, preserving all experts but making each one compact.

\paragraph{Elastic and dynamic architectures.}
Slimmable Networks~\citep{yu2019universally} train a single network that can execute at multiple widths. Once-for-All~\citep{cai2020once} extends this to depth, width, and resolution jointly. DynaBERT~\citep{hou2020dynabert} applies elastic width and depth to transformers. \citet{chen2024dynamic} provide a comprehensive survey of dynamic neural networks, including input-adaptive depth, width, and routing decisions. These approaches typically use knowledge distillation across sub-networks. KaleidoNet differs by using learnable per-neuron Gumbel-sigmoid gates that converge to hard decisions during training, guided by the cubic schedule rather than distillation targets.

\paragraph{Gumbel-based discrete optimization.}
The Gumbel-Softmax trick~\citep{jang2017categorical,maddison2017concrete} enables gradient-based optimization of discrete variables via continuous relaxation. Straight-through estimators~\citep{bengio2013estimating,courbariaux2016binarized} pass gradients through hard thresholds by using the soft gradient in the backward pass. We use a Gumbel-sigmoid variant with straight-through estimation for binary per-neuron masks, annealing the temperature from $\tau{=}5.0$ to $\tau{=}0.1$ over training.

\paragraph{Positioning.}
Prior ViT pruning methods~\citep{chen2021chasing,yu2022width} operate on standard (non-MoE) architectures where Lagrangian or regularization-based approaches work well (e.g., 68.4\% on ImageNet-1K with pruned DeiT-S at 30\% sparsity). Our work targets a fundamentally different setting---pruning \emph{within} MoE expert sub-networks---where we show that these same Lagrangian approaches catastrophically fail due to gradient-scale mismatch. KaleidoNet's cubic schedule is specifically designed for this MoE-internal pruning regime.

\paragraph{Concurrent work.}
Several concurrent works address MoE compression from complementary angles. SlimMoE~\citep{slimmoe2025} retains all experts and prunes redundant neurons within each via iterative prune-and-distill. STUN~\citep{stun2025} combines structured (expert-level) and unstructured (weight-level) pruning for scalable MoE compression. MoNE~\citep{mone2025} replaces redundant experts with lightweight novice networks. MoE-Pruner~\citep{moepruner2025} uses router-guided magnitude pruning at the neuron level. None of these works analyze the Lagrangian failure mode; our contribution is the diagnosis of \emph{why} Lagrangian approaches fail at per-neuron granularity, which informs the design of future MoE compression methods regardless of the specific pruning strategy adopted.
